% GENERATED FILE. Edit canonical lessons, then run npm run generate:books.
% canonical-source-hash: fnv1a64:cfd598faafdf0e52
% canonical-lessons: JA-C285-okonau, JA-C285-hiraku, JA-C285-kuwawaru, JA-C285-kinen, JA-C285-shukujitsu, JA-R285-review-280-285

\chapter{A Big Day in Town}
\label{ch:ja-a-big-day-in-town}
\hlchaptermodality{285}

\begin{chapteropening}
\textbf{By the end of this chapter:} \emph{I can say what the town held and who joined in.}

Complete the last lesson of chapter 285: I can say what the town held and who joined in.
\end{chapteropening}

\section[\texorpdfstring{\ja{おこなう} (okonau)}{おこなう (okonau)}]{\ja{おこなう} (okonau) — to carry out; to hold (an event)}

\label{lesson:JA-C285-okonau}



\begin{quote}
Before the new one: say the Japanese for to compete, then the Japanese for bad.
\end{quote}

\subsection*{You'll want to know: \ja{おこなう}}
\textbf{\ja{おこなう}} — \emph{okonau} — \textquotedblleft{}to carry out; to hold (an event)\textquotedblright{}.

\begin{quote}
\textbf{\ja{らいしゅう} \ja{まちで} \ja{まつりを} \ja{おこないます。}} — \emph{raishū machi de matsuri o okonaimasu} — \textquotedblleft{}The town holds a festival next week.\textquotedblright{}
\end{quote}

Polite \textbf{\ja{おこないます}}, \emph{okonaimasu}; a godan verb. Past (polite): \textbf{\ja{がっこうが} \ja{しけんを} \ja{おこないました。}}, \emph{gakkō ga shiken o okonaimashita}.

\begin{cousinweb}
\textbf{\ja{おこなう}} (\emph{okonau}) is a native Japanese word (wago), inherited rather than borrowed from Chinese; its kanji was attached later as a kun'yomi (JMdict).
\end{cousinweb}

\begin{grammarlens}[title={What you've built}]
One more word for telling what happened.
\end{grammarlens}

\subsection*{Your turn}
\begin{itemize}
\raggedright
  \item \emph{Say it:} \emph{okonau}
  \item \emph{Say it:} \emph{okonau}, once more
  \item \emph{Say it:} \emph{raishū machi de matsuri o okonaimasu} — The town holds a festival next week.
  \item \emph{Recall:} say the Japanese for to compete, then the Japanese for bad, then say \emph{okonau} again
\end{itemize}

\begin{tcolorbox}[breakable,colback=teal!4,colframe=teal!35!black,title={Before you move on}]
What does \ja{おこなう} mean? (\textquotedblleft{}To carry out; to hold (an event)\textquotedblright{}.) Say it once more.
\end{tcolorbox}
\section[\texorpdfstring{\ja{ひらく} (hiraku)}{ひらく (hiraku)}]{\ja{ひらく} (hiraku) — to hold (an event); to open}

\label{lesson:JA-C285-hiraku}



\begin{quote}
Before the new one: say the Japanese for bad, then the Japanese for to carry out.
\end{quote}

\subsection*{You'll want to know: \ja{ひらく}}
\textbf{\ja{ひらく}} — \emph{hiraku} — \textquotedblleft{}to hold (an event); to open\textquotedblright{}.

\begin{quote}
\textbf{\ja{にちようびに} \ja{こうえんで} \ja{まつりを} \ja{ひらきます。}} — \emph{nichiyōbi ni kōen de matsuri o hirakimasu} — \textquotedblleft{}We hold a festival in the park on Sunday.\textquotedblright{}
\end{quote}

Polite \textbf{\ja{ひらきます}}, \emph{hirakimasu}; a godan verb. Past (polite): \textbf{\ja{せいふが} \ja{かいぎを} \ja{ひらきました。}}, \emph{seifu ga kaigi o hirakimashita}.

\begin{cousinweb}
\textbf{\ja{ひらく}} (\emph{hiraku}) is a native Japanese word (wago), inherited rather than borrowed from Chinese; its kanji was attached later as a kun'yomi (JMdict).
\end{cousinweb}

\begin{grammarlens}[title={What you've built}]
One more word for telling what happened.
\end{grammarlens}

\subsection*{Your turn}
\begin{itemize}
\raggedright
  \item \emph{Say it:} \emph{hiraku}
  \item \emph{Say it:} \emph{hiraku}, once more
  \item \emph{Say it:} \emph{nichiyōbi ni kōen de matsuri o hirakimasu} — We hold a festival in the park on Sunday.
  \item \emph{Recall:} say the Japanese for bad, then the Japanese for to carry out, then say \emph{hiraku} again
\end{itemize}

\begin{tcolorbox}[breakable,colback=teal!4,colframe=teal!35!black,title={Before you move on}]
What does \ja{ひらく} mean? (\textquotedblleft{}To hold (an event); to open\textquotedblright{}.) Say it once more.
\end{tcolorbox}
\section[\texorpdfstring{\ja{くわわる} (kuwawaru)}{くわわる (kuwawaru)}]{\ja{くわわる} (kuwawaru) — to join (in); to take part}

\label{lesson:JA-C285-kuwawaru}



\begin{quote}
Before the new one: say the Japanese for to carry out, then the Japanese for to hold (an event).
\end{quote}

\subsection*{You'll want to know: \ja{くわわる}}
\textbf{\ja{くわわる}} — \emph{kuwawaru} — \textquotedblleft{}to join (in); to take part\textquotedblright{}.

\begin{quote}
\textbf{\ja{らいねんは} \ja{ちちも} \ja{まつりに} \ja{くわわるでしょう。}} — \emph{rainen wa chichi mo matsuri ni kuwawaru deshō} — \textquotedblleft{}Next year my father will probably join the festival too.\textquotedblright{}
\end{quote}

Polite \textbf{\ja{くわわります}}, \emph{kuwawarimasu}; a godan verb. Past (polite): \textbf{\ja{ともだちが} \ja{くわわりました。}}, \emph{tomodachi ga kuwawarimashita}.

\begin{cousinweb}
\textbf{\ja{くわわる}} (\emph{kuwawaru}) is a native Japanese word (wago), inherited rather than borrowed from Chinese; its kanji was attached later as a kun'yomi (JMdict).
\end{cousinweb}

\begin{grammarlens}[title={What you've built}]
One more word for telling what happened.
\end{grammarlens}

\subsection*{Your turn}
\begin{itemize}
\raggedright
  \item \emph{Say it:} \emph{kuwawaru}
  \item \emph{Say it:} \emph{kuwawaru}, once more
  \item \emph{Say it:} \emph{rainen wa chichi mo matsuri ni kuwawaru deshō} — Next year my father will probably join the festival too.
  \item \emph{Recall:} say the Japanese for to carry out, then the Japanese for to hold (an event), then say \emph{kuwawaru} again
\end{itemize}

\begin{tcolorbox}[breakable,colback=teal!4,colframe=teal!35!black,title={Before you move on}]
What does \ja{くわわる} mean? (\textquotedblleft{}To join (in); to take part\textquotedblright{}.) Say it once more.
\end{tcolorbox}
\section[\texorpdfstring{\ja{きねん} (kinen)}{きねん (kinen)}]{\ja{きねん} (kinen) — a commemoration; a memento}

\label{lesson:JA-C285-kinen}



\begin{quote}
Before the new one: say the Japanese for to hold (an event), then the Japanese for to join (in).
\end{quote}

\subsection*{You'll want to know: \ja{きねん}}
\textbf{\ja{きねん}} — \emph{kinen} — \textquotedblleft{}a commemoration; a memento\textquotedblright{}.

\begin{quote}
\textbf{\ja{きねんに} \ja{しゃしんを} \ja{とりました。}} — \emph{kinen ni shashin o torimashita} — \textquotedblleft{}We took a photo to remember the day.\textquotedblright{}
\end{quote}

\begin{cousinweb}
\textbf{\ja{きねん}} (\emph{kinen}) is Sino-Japanese (kango): two Chinese characters read with their on'yomi, readings borrowed from Chinese. They are \emph{ki}, \textquotedblleft{}account\textquotedblright{} (Middle Chinese \emph{kiH}) and \emph{nen}, \textquotedblleft{}thought\textquotedblright{} (\emph{nemH}) (KANJIDIC2; Middle Chinese in Baxter's transcription).
\end{cousinweb}

\begin{grammarlens}[title={What you've built}]
One more word for telling what happened.
\end{grammarlens}

\subsection*{Your turn}
\begin{itemize}
\raggedright
  \item \emph{Say it:} \emph{kinen}
  \item \emph{Say it:} \emph{kinen}, once more
  \item \emph{Say it:} \emph{kinen ni shashin o torimashita} — We took a photo to remember the day.
  \item \emph{Recall:} say the Japanese for to hold (an event), then the Japanese for to join (in), then say \emph{kinen} again
\end{itemize}

\begin{tcolorbox}[breakable,colback=teal!4,colframe=teal!35!black,title={Before you move on}]
What does \ja{きねん} mean? (\textquotedblleft{}A commemoration; a memento\textquotedblright{}.) Say it once more.
\end{tcolorbox}
\section[\texorpdfstring{\ja{しゅくじつ} (shukujitsu)}{しゅくじつ (shukujitsu)}]{\ja{しゅくじつ} (shukujitsu) — a public holiday}

\label{lesson:JA-C285-shukujitsu}



\begin{quote}
Before the new one: say the Japanese for to join (in), then the Japanese for a commemoration.
\end{quote}

\subsection*{You'll want to know: \ja{しゅくじつ}}
\textbf{\ja{しゅくじつ}} — \emph{shukujitsu} — \textquotedblleft{}a public holiday\textquotedblright{}.

\begin{quote}
\textbf{\ja{あしたは} \ja{しゅくじつです。}} — \emph{ashita wa shukujitsu desu} — \textquotedblleft{}Tomorrow is a public holiday.\textquotedblright{}
\end{quote}

\begin{cousinweb}
\textbf{\ja{しゅくじつ}} (\emph{shukujitsu}) is Sino-Japanese (kango): two Chinese characters read with their on'yomi, readings borrowed from Chinese. They are \emph{shuku}, \textquotedblleft{}celebrate\textquotedblright{} (Middle Chinese \emph{tsyuwk}) and \emph{jitsu}, \textquotedblleft{}day\textquotedblright{} (\emph{nyit}) (KANJIDIC2; Middle Chinese in Baxter's transcription).
\end{cousinweb}

\begin{grammarlens}[title={What you've built}]
One more word for telling what happened.
\end{grammarlens}

\subsection*{Your turn}
\begin{itemize}
\raggedright
  \item \emph{Say it:} \emph{shukujitsu}
  \item \emph{Say it:} \emph{shukujitsu}, once more
  \item \emph{Say it:} \emph{ashita wa shukujitsu desu} — Tomorrow is a public holiday.
  \item \emph{Recall:} say the Japanese for to join (in), then the Japanese for a commemoration, then say \emph{shukujitsu} again
\end{itemize}

\begin{tcolorbox}[breakable,colback=teal!4,colframe=teal!35!black,title={Before you move on}]
What does \ja{しゅくじつ} mean? (\textquotedblleft{}A public holiday\textquotedblright{}.) Say it once more.
\end{tcolorbox}
\section[(dialogue)]{Looking back: from What the News Said to A Big Day in Town}

\label{lesson:JA-R285-review-280-285}



\begin{quote}
Say the Japanese for a newspaper article and to print (in a paper).
\end{quote}

\subsection*{Your turn}
Say these aloud:

\begin{itemize}
\raggedright
  \item a newspaper article, to print (in a paper), an event, a reporter, to hand out
  \item a rumour, to be passed on, to spread widely, a topic of conversation, an uproar
  \item the world, a war, peace, to fight, the government
  \item to announce, to report, to investigate, information, a fact
  \item an election, the people (of a country), the prime minister, to compete, bad
  \item to carry out, to hold (an event), to join (in), a commemoration, a public holiday
\end{itemize}

\begin{tcolorbox}[breakable,colback=teal!4,colframe=teal!35!black,title={Before you move on}]
A newspaper article? (\textbf{\ja{きじ}}, \emph{kiji}.) To print (in a paper)? (\textbf{\ja{のせる}}, \emph{noseru}.) An event? (\textbf{\ja{できごと}}, \emph{dekigoto}.) A reporter? (\textbf{\ja{きしゃ}}, \emph{kisha}.) To hand out? (\textbf{\ja{くばる}}, \emph{kubaru}.) A rumour? (\textbf{\ja{うわさ}}, \emph{uwasa}.) To be passed on? (\textbf{\ja{つたわる}}, \emph{tsutawaru}.) To spread widely? (\textbf{\ja{ひろまる}}, \emph{hiromaru}.) A topic of conversation? (\textbf{\ja{わだい}}, \emph{wadai}.) An uproar? (\textbf{\ja{おおさわぎ}}, \emph{ōsawagi}.) The world? (\textbf{\ja{せかい}}, \emph{sekai}.) A war? (\textbf{\ja{せんそう}}, \emph{sensō}.) Peace? (\textbf{\ja{へいわ}}, \emph{heiwa}.) To fight? (\textbf{\ja{たたかう}}, \emph{tatakau}.) The government? (\textbf{\ja{せいふ}}, \emph{seifu}.) To announce? (\textbf{\ja{はっぴょうする}}, \emph{happyō suru}.) To report? (\textbf{\ja{ほうこくする}}, \emph{hōkoku suru}.) To investigate? (\textbf{\ja{ちょうさする}}, \emph{chōsa suru}.) Information? (\textbf{\ja{じょうほう}}, \emph{jōhō}.) A fact? (\textbf{\ja{じじつ}}, \emph{jijitsu}.) An election? (\textbf{\ja{せんきょ}}, \emph{senkyo}.) The people (of a country)? (\textbf{\ja{こくみん}}, \emph{kokumin}.) The prime minister? (\textbf{\ja{しゅしょう}}, \emph{shushō}.) To compete? (\textbf{\ja{あらそう}}, \emph{arasou}.) Bad? (\textbf{\ja{わるい}}, \emph{warui}.) To carry out? (\textbf{\ja{おこなう}}, \emph{okonau}.) To hold (an event)? (\textbf{\ja{ひらく}}, \emph{hiraku}.) To join (in)? (\textbf{\ja{くわわる}}, \emph{kuwawaru}.) A commemoration? (\textbf{\ja{きねん}}, \emph{kinen}.) A public holiday? (\textbf{\ja{しゅくじつ}}, \emph{shukujitsu}.)
\end{tcolorbox}
